\begin{table*}[t]
\centering
\caption{R-AUROC for downstream linear-probe classification with the frozen encoder (0\% fine-tuning). Each entry is the mean $\pm$ standard deviation across evaluation units.}
\label{tab:har_linear_probe_0}
\small
\setlength{\tabcolsep}{9pt}
\renewcommand{\arraystretch}{1.17}
\begin{tabular}{@{}lcccc@{}}
\toprule
Dataset & CE & Gaussian & Student-$t$ & MC Dropout \\
\midrule
CAPTURE-24 & $0.416 \pm 0.089$ & $0.639 \pm 0.071$ & $0.655 \pm 0.083$ & $0.421 \pm 0.105$ \\
PAMAP2 & $0.482 \pm 0.084$ & $0.359 \pm 0.081$ & $0.297 \pm 0.090$ & $0.336 \pm 0.081$ \\
RealWorld & $0.624 \pm 0.103$ & $0.365 \pm 0.121$ & $0.348 \pm 0.106$ & $0.354 \pm 0.075$ \\
USC-HAD & $0.344 \pm 0.074$ & $0.576 \pm 0.086$ & $0.540 \pm 0.072$ & $0.545 \pm 0.091$ \\
HARTH & $0.431 \pm 0.097$ & $0.483 \pm 0.239$ & $0.462 \pm 0.228$ & $0.303 \pm 0.122$ \\
\bottomrule
\end{tabular}
\par\vspace{3pt}
\begin{minipage}{0.94\textwidth}
\footnotesize\emph{Note.} R-AUROC $>0.5$ denotes conventional error ranking; R-AUROC $<0.5$ denotes inverse ranking. Values are averaged per held-out unit (CAPTURE-24: 15 ten-subject test folds; other datasets: LOSO subjects). These directions are descriptive, not significance claims.
\end{minipage}
\end{table*}
